%# =====================================================================================================
%# ptreport default LaTeX template - "Corporate Band"
%#
%# Jinja with LaTeX-friendly delimiters:
%#   \VAR{value}          prints a value; it is ALWAYS escaped for LaTeX, so report content is safe
%#   \BLOCK{if ...}       statements: if / for / set ... closed with \BLOCK{endif}, \BLOCK{endfor}
%#   lines starting with %# are template comments (removed from the output)
%# Filters: richtext, verbatim, verbatim_file, url, url_cwe, reference_list, image, asset, score, breakable, filesize,
%#   cvss_rating_key, cvss_rating_label, vector_strip, cwe_name, tex.
%# Data: report, brand (color, logo), versions, targets, findings, counts, severities, L (labels),
%#   lang, cvss_header, show_finding_history, generated_at. Each finding has files (filename, description,
%#   size, sha256, mime, uploaded_at).
%#
%# Download it from Settings -> PDF template, edit, upload it back. Start with the EDIT HERE block.
%# =====================================================================================================
\documentclass[10pt,a4paper]{article}

% ======================================================================================== EDIT HERE ==
% Brand colour (Settings -> Branding). Replace with a fixed value such as {1D3557} if you prefer.
\newcommand{\ptBrandHex}{\VAR{brand.color}}
% Fonts: TeX Gyre Heros (Helvetica-like) for text, DejaVu Sans Mono for data. Any OTF/TTF works.
\newcommand{\ptSetFonts}{%
  \setmainfont{texgyreheros}[Extension=.otf, UprightFont=*-regular, BoldFont=*-bold, ItalicFont=*-italic, BoldItalicFont=*-bolditalic]%
  \setmonofont{DejaVuSansMono}[Extension=.ttf, BoldFont=*-Bold, Scale=0.85]}
% Logo (Settings -> Branding). Empty = the company name is printed instead.
\def\ptLogoFile{\VAR{brand.logo|asset('logo')}}
% Confidentiality label on the cover and in every footer (report field "Confidentiality label").
\def\ptClassification{\VAR{report.classification}}
% Severity colours (also used by the CVSS rating chips).
\newcommand{\ptSeverityColors}{%
  \definecolor{sevcritical}{HTML}{9E1C34}\definecolor{sevhigh}{HTML}{C4470F}\definecolor{sevmedium}{HTML}{B8860B}%
  \definecolor{sevlow}{HTML}{2E7D4F}\definecolor{sevinfo}{HTML}{2F6DB5}\definecolor{sevnone}{HTML}{6B7680}}
% ===================================================================================================

\usepackage{fontspec}
\ptSetFonts
\usepackage[\BLOCK{if lang == 'pl'}polish\BLOCK{else}english\BLOCK{endif}]{babel}
\usepackage[a4paper,top=26mm,bottom=24mm,left=20mm,right=20mm,headheight=16pt,headsep=8mm,footskip=11mm]{geometry}
\usepackage[table]{xcolor}
\usepackage{graphicx,array,tabularx,longtable,booktabs,fancyhdr,lastpage,titlesec,titletoc,enumitem,caption,needspace}
\usepackage{tikz}
\usepackage[most]{tcolorbox}
\usepackage{fvextra}
\usepackage[hidelinks,unicode]{hyperref}
\usepackage{xurl}

% ------------------------------------------------------------------------------------------ colours --
\definecolor{brand}{HTML}{\ptBrandHex}
\colorlet{brandsoft}{brand!7!white}
\colorlet{brandline}{brand!35!white}
\definecolor{ink}{HTML}{1A2128}
\definecolor{muted}{HTML}{5A6570}
\definecolor{line}{HTML}{D9DEE3}
\definecolor{soft}{HTML}{F4F6F8}
\ptSeverityColors
\colorlet{tone0}{soft}\colorlet{tone1}{line}\colorlet{tone2}{muted!55}\colorlet{tone3}{ink!85}

% -------------------------------------------------------------------------------------- typography --
\color{ink}
\setlength{\parindent}{0pt}
\setlength{\parskip}{0.6em}
\linespread{1.1}
\setlist{itemsep=0.2em,topsep=0.3em,leftmargin=1.4em}
\captionsetup{font={small,color=muted},labelfont={bf,color=brand},skip=6pt}
\renewcommand{\arraystretch}{1.45}
\newcolumntype{L}{>{\raggedright\arraybackslash}X}

\titleformat{\section}{\LARGE\bfseries\color{ink}}{\textcolor{brand}{\thesection}}{0.7em}{}[\vspace{2pt}{\color{brandline}\titlerule[0.8pt]}]
\titleformat{\subsection}{\large\bfseries\color{ink}}{\textcolor{brand}{\thesubsection}}{0.6em}{}
\titleformat{\subsubsection}{\normalsize\bfseries\color{ink}}{\textcolor{brand}{\thesubsubsection}}{0.6em}{}
\titlespacing*{\section}{0pt}{0.4em}{1.1em}
\titlespacing*{\subsection}{0pt}{1.2em}{0.4em}
\titlespacing*{\subsubsection}{0pt}{1em}{0.3em}
\setcounter{secnumdepth}{3}
\setcounter{tocdepth}{3}
% Labels of the blocks inside a finding (Description, Impact ...): not part of the numbering.
\newcommand{\ptBlockTitle}[1]{\par\needspace{8\baselineskip}\vspace{1.1em}{\small\bfseries\color{brand}\MakeUppercase{#1}}\par\nopagebreak\vspace{0.3em}}
\titlecontents{section}[0pt]{\addvspace{0.7em}\bfseries}{\textcolor{brand}{\thecontentslabel}\quad}{}{\hfill\contentspage}
\titlecontents{subsection}[4.3em]{\addvspace{0.1em}\small}{\contentslabel{2.4em}}{\hspace*{-2.4em}}{\titlerule*[0.6pc]{.}\contentspage}
\titlecontents{subsubsection}[7.3em]{\small\color{muted}}{\contentslabel{3em}}{\hspace*{-3em}}{\titlerule*[0.6pc]{.}\contentspage}

% -------------------------------------------------------------------------------- header / footer --
\pagestyle{fancy}
\fancyhf{}
\fancyhead[L]{\ifx\ptLogoFile\empty{\small\bfseries\color{brand}\VAR{report.vendor_name}}\else\raisebox{-0.3\height}{\includegraphics[height=6mm]{\ptLogoFile}}\fi}
\fancyhead[R]{\small\color{muted}\VAR{report.title}}
\fancyfoot[L]{\ifx\ptClassification\empty\else{\footnotesize\bfseries\color{brand}\ptClassification}\fi}
\BLOCK{set year = (report.report_date or report.end_date or report.start_date or generated_at)[:4]}
\fancyfoot[C]{\footnotesize\color{muted}\BLOCK{if report.vendor_name}\textcopyright{} \VAR{year} \VAR{report.vendor_name}\BLOCK{endif}\BLOCK{if report.report_version} · \VAR{L.version} \VAR{report.report_version}\BLOCK{endif}}
\fancyfoot[R]{\footnotesize\color{muted}\thepage\,/\,\pageref*{LastPage}}
\renewcommand{\headrule}{{\color{brandline}\hrule height 0.6pt}}
\renewcommand{\footrulewidth}{0pt}

% ---------------------------------------------------------------------------------------- elements --
% Severity: the tester's assessment, a filled badge.
\newtcbox{\ptSevBox}[1]{on line, colback=sev#1, colframe=sev#1, coltext=white, boxrule=0pt, arc=2pt,
  boxsep=0pt, left=4pt, right=4pt, top=2pt, bottom=2pt, fontupper=\scriptsize\bfseries}
\newcommand{\ptSeverity}[2]{\ptSevBox{#1}{\MakeUppercase{#2}}}
% Table column: a share of the line width, net of the column gaps (so tables never run into the margin).
\newcolumntype{W}[1]{p{\dimexpr #1\linewidth-2\tabcolsep\relax}}
% Retest result: a coloured dot and the result (fixed / partially fixed / not fixed / not tested / none).
\colorlet{rtfixed}{sevlow}\colorlet{rtpartiallyfixed}{sevmedium}\colorlet{rtnotfixed}{sevcritical}
\colorlet{rtnottested}{muted}\colorlet{rtnone}{muted}
\newcommand{\ptRetest}[2]{{\color{rt#1}\raisebox{0.15ex}{\tikz\fill (0,0) circle (0.28em);}\,\bfseries\small #2}}
% CVSS: the calculated score coloured by its own rating; outlined so it never reads as the severity badge.
\newtcbox{\ptCvssBox}[1]{on line, colback=white, colframe=sev#1, coltext=sev#1, boxrule=0.8pt, arc=2pt,
  boxsep=0pt, left=4pt, right=4pt, top=1.5pt, bottom=1.5pt, fontupper=\scriptsize\bfseries}
\newcommand{\ptCvss}[3]{\ptCvssBox{#1}{\texttt{#2}\enspace #3}}
% Risk band on top of each finding's facts: the severity in large type on its own colour.
\newcommand{\ptRiskBand}[2]{{\Large\bfseries\MakeUppercase{#1}}\hfill{\normalsize\mdseries #2}}
% One cell of the CVSS vector strip; tone 0-3 = how much the value raises the score.
\newcommand{\ptCell}[3]{{\setlength{\fboxsep}{1.5pt}\colorbox{tone#1}{\textcolor{\ifnum#1>1 white\else ink\fi}{\scriptsize\texttt{#2:#3}}}}\hspace{1pt}}
\newcommand{\ptUrl}[1]{\textcolor{brand}{#1}}
% Raw text (HTTP messages, code) is read from files and never interpreted as LaTeX.
\newtcolorbox{ptCodeBox}[1][]{enhanced, breakable, colback=soft, colframe=line, boxrule=0.5pt, arc=2pt,
  left=6pt, right=6pt, top=5pt, bottom=5pt, before skip=4pt, after skip=8pt, #1}
\newcommand{\ptVerbatim}[1]{\begin{ptCodeBox}\VerbatimInput[breaklines,breakanywhere,fontsize=\footnotesize,tabsize=4]{#1}\end{ptCodeBox}}
\newcommand{\ptHttp}[2]{\begin{ptCodeBox}[title={#1}, attach boxed title to top left={xshift=6pt,yshift=-2.5mm},
  boxed title style={colback=brand, colframe=brand, arc=1.5pt, boxrule=0pt, top=0.5pt, bottom=0.5pt},
  fonttitle=\scriptsize\bfseries\color{white}, top=9pt]\VerbatimInput[breaklines,breakanywhere,fontsize=\footnotesize,tabsize=4]{#2}\end{ptCodeBox}}
\newcommand{\ptImage}[1]{\fcolorbox{line}{white}{\includegraphics[width=\dimexpr\linewidth-2\fboxsep-2\fboxrule\relax,height=0.5\textheight,keepaspectratio]{#1}}}
\newcommand{\ptMissingImage}{\fbox{\parbox{0.9\linewidth}{\centering\small\color{muted}(image format not supported in PDF)}}}
\newcommand{\ptLabel}[1]{{\scriptsize\bfseries\color{muted}\MakeUppercase{#1}}}
\newcommand{\ptHead}[1]{{\color{white}\scriptsize\bfseries\MakeUppercase{#1}}}
% Tables written in report text (Markdown "| a | b |"): zebra rows, a brand-coloured header row.
\newcommand{\ptTableBegin}{\par\rowcolors{2}{soft}{white}}
\newcommand{\ptTableEnd}{\rowcolors{1}{}{}\par}
\newcommand{\ptTableHead}[1]{\cellcolor{brand}{\color{white}\scriptsize\bfseries #1}}

\hypersetup{pdftitle={\VAR{report.title}}, pdfauthor={\VAR{report.vendor_name or report.assessor}}}

\begin{document}

%# --------------------------------------------------------------------------------------------- cover
\begin{titlepage}
\thispagestyle{empty}
\begin{tikzpicture}[remember picture, overlay]
  \fill[brand] (current page.north west) rectangle ([yshift=-118mm]current page.north east);
  \fill[brand!75!black] ([yshift=-118mm]current page.north west) rectangle ([yshift=-120.5mm]current page.north east);
  \node[anchor=north west, inner sep=0pt] at ([xshift=20mm,yshift=-18mm]current page.north west)
    {\ifx\ptLogoFile\empty{\color{white}\Large\bfseries \VAR{report.vendor_name}}\else{\setlength{\fboxsep}{3mm}\colorbox{white}{\includegraphics[height=12mm]{\ptLogoFile}}}\fi};
  \node[anchor=south west, text width=170mm, align=left, inner sep=0pt] at ([xshift=20mm,yshift=-106mm]current page.north west)
    {{\color{white!78!brand}\small\bfseries\MakeUppercase{\VAR{L.report}}}\\[5mm]
     {\color{white}\fontsize{26}{31}\selectfont\bfseries\spaceskip=0.34em plus 0.12em minus 0.06em\relax\VAR{report.title}}\BLOCK{if report.client_name}\\[6mm]
     {\color{white!85!brand}\large \VAR{report.client_name}}\BLOCK{endif}};
\end{tikzpicture}
\vspace*{110mm}

\begin{tabularx}{\linewidth}{@{}>{\leavevmode\color{muted}\small\bfseries}p{40mm}L@{}}
\BLOCK{if report.client_name}\VAR{L.client} & \VAR{report.client_name} \\ \BLOCK{endif}
\BLOCK{if report.assessor}\VAR{L.assessor} & \VAR{report.assessor} \\ \BLOCK{endif}
\BLOCK{if report.start_date or report.end_date}\VAR{L.period} & \VAR{report.start_date or ''}\BLOCK{if report.start_date and report.end_date} -- \BLOCK{endif}\VAR{report.end_date or ''} \\ \BLOCK{endif}
\BLOCK{if report.vendor_name}\VAR{L.prepared_by} & \VAR{report.vendor_name} \\ \BLOCK{endif}
\BLOCK{if report.vendor_contact}\VAR{L.contact} & \VAR{report.vendor_contact} \\ \BLOCK{endif}
\BLOCK{if report.report_version}\VAR{L.version} & \VAR{report.report_version} \\ \BLOCK{endif}
\BLOCK{if report.report_date}\VAR{L.report_date} & \VAR{report.report_date} \\ \BLOCK{endif}
\end{tabularx}

\vfill
\ifx\ptClassification\empty\else
\tcbox[on line, colback=white, colframe=brand, coltext=brand, boxrule=1pt, arc=2pt, left=8pt, right=8pt, top=4pt, bottom=4pt, fontupper=\bfseries]{\ptClassification}\par\vspace{4mm}
\fi
{\footnotesize\color{muted}\VAR{L.generated}: \VAR{generated_at}}
\end{titlepage}

%# ------------------------------------------------------------------------------- version history, toc
\BLOCK{if versions}
\section*{\VAR{L.version_history}}
\rowcolors{2}{soft}{white}
\begin{tabularx}{\linewidth}{@{}p{18mm}p{24mm}\BLOCK{if versions_have_author}p{40mm}\BLOCK{endif}L@{}}
\rowcolor{white}\ptLabel{\VAR{L.version}} & \ptLabel{\VAR{L.date}} & \BLOCK{if versions_have_author}\ptLabel{\VAR{L.author}} & \BLOCK{endif}\ptLabel{\VAR{L.description}} \\
\BLOCK{for v in versions}
\texttt{\VAR{v.version}} & \texttt{\VAR{v.date or ''}} & \BLOCK{if versions_have_author}\VAR{v.author} & \BLOCK{endif}\VAR{v.description} \\
\BLOCK{endfor}
\end{tabularx}
\rowcolors{1}{}{}
\vspace{1.5em}
\BLOCK{endif}

\renewcommand{\contentsname}{\VAR{L.toc}}
\tableofcontents
\BLOCK{if fingerprint}
\vspace{1.2em}
{\scriptsize\color{muted}\VAR{L.fingerprint}: \texttt{\VAR{fingerprint}}\par}
\BLOCK{endif}
\newpage

%# ================================================================================ building blocks ==
%# Each block below is a Jinja macro; the report body at the end decides their order (report structure).

\BLOCK{macro heading(sec)}
\BLOCK{if sec.new_page}\clearpage\BLOCK{endif}
\BLOCK{if sec.level == 1}
\BLOCK{if sec.numbered}\section{\VAR{sec.title}}\BLOCK{else}\section*{\VAR{sec.title}}\addcontentsline{toc}{section}{\VAR{sec.title}}\BLOCK{endif}
\BLOCK{else}
\BLOCK{if sec.numbered}\subsection{\VAR{sec.title}}\BLOCK{else}\subsection*{\VAR{sec.title}}\addcontentsline{toc}{subsection}{\VAR{sec.title}}\BLOCK{endif}
\BLOCK{endif}
\BLOCK{endmacro}

\BLOCK{macro targets_table()}
\rowcolors{2}{soft}{white}
\begin{tabularx}{\linewidth}{@{}p{45mm}p{32mm}L@{}}
\rowcolor{white}\ptLabel{\VAR{L.target}} & \ptLabel{\VAR{L.type}} & \ptLabel{\VAR{L.address}} \\
\BLOCK{for t in targets}
\textbf{\VAR{t.name}} & \VAR{TARGET_KIND_LABELS[lang][t.kind]} & \texttt{\VAR{t.address}} \\
\BLOCK{endfor}
\end{tabularx}
\rowcolors{1}{}{}
\BLOCK{endmacro}

\BLOCK{macro severity_bars()}
\BLOCK{set maxcount = [counts.critical, counts.high, counts.medium, counts.low, counts.info, 1]|max}
\begin{tcolorbox}[enhanced, colback=white, colframe=line, boxrule=0.5pt, arc=3pt, left=10pt, right=10pt, top=10pt, bottom=8pt]
\centering
\begin{tikzpicture}[x=1mm, y=1mm]
\BLOCK{for s in severities}
  \node[anchor=east, font=\small] at (28, \VAR{(-8 * loop.index0)|round(2)}) {\VAR{sev_labels[s]}};
  \fill[soft, rounded corners=1pt] (31, \VAR{(-8 * loop.index0 - 2.2)|round(2)}) rectangle (131, \VAR{(-8 * loop.index0 + 2.2)|round(2)});
\BLOCK{if counts[s]}
  \fill[sev\VAR{s}, rounded corners=1pt] (31, \VAR{(-8 * loop.index0 - 2.2)|round(2)}) rectangle (\VAR{(31 + 100 * counts[s] / maxcount)|round(2)}, \VAR{(-8 * loop.index0 + 2.2)|round(2)});
\BLOCK{endif}
  \node[anchor=west, font=\small\bfseries] at (134, \VAR{(-8 * loop.index0)|round(2)}) {\VAR{counts[s]}};
\BLOCK{endfor}
  \draw[line] (0, -38) -- (145, -38);
  \node[anchor=east, font=\small\bfseries] at (28, -43) {\VAR{L.total}};
  \node[anchor=west, font=\small\bfseries] at (134, -43) {\VAR{findings|length}};
\end{tikzpicture}
\end{tcolorbox}
\BLOCK{endmacro}

\BLOCK{macro findings_table()}
\BLOCK{if findings}
\rowcolors{2}{soft}{white}
\begin{longtable}{@{}W{.06}W{.13}W{.36}W{.14}W{.18}>{\raggedright\arraybackslash}W{.13}@{}}
\rowcolor{brand}\ptHead{\VAR{L.no}} & \ptHead{\VAR{L.id}} & \ptHead{\VAR{L.title}} & \ptHead{\VAR{L.severity}} & \ptHead{\VAR{cvss_header}} & \ptHead{\VAR{L.status}} \\
\endhead
\BLOCK{for f in findings}
\small\VAR{'%02d.'|format(loop.index)} & \hyperref[f:\VAR{f.id}]{\small\texttt{\VAR{f.display_id}}} & \raggedright\small \VAR{f.title} & \ptSeverity{\VAR{f.severity}}{\VAR{sev_labels[f.severity]}} & \BLOCK{if f.cvss_score is not none}\ptCvss{\VAR{f|cvss_rating_key}}{\VAR{f.cvss_score|score}}{\VAR{f|cvss_rating_label(lang)}}\BLOCK{else}{\color{muted}--}\BLOCK{endif} & \small\VAR{status_labels[f.status]} \\
\BLOCK{endfor}
\end{longtable}
\rowcolors{1}{}{}
\BLOCK{else}
\VAR{L.no_findings}
\BLOCK{endif}
\BLOCK{endmacro}

\BLOCK{macro retest_summary()}
\BLOCK{if original_report}
\VAR{L.original_report}: \textbf{\VAR{original_report.title}}\BLOCK{if original_report.report_date} (\VAR{L.original_date} \texttt{\VAR{original_report.report_date}})\BLOCK{endif}\par\medskip
\BLOCK{endif}
\BLOCK{if findings}
\noindent\BLOCK{for key, label in retest_labels.items()}\BLOCK{if retest_counts[key]}\ptRetest{\VAR{(key or 'none')|replace('_', '')}}{\VAR{label}}: \textbf{\VAR{retest_counts[key]}}\qquad \BLOCK{endif}\BLOCK{endfor}\par\medskip
\rowcolors{2}{soft}{white}
\begin{longtable}{@{}W{.10}W{.33}W{.17}>{\raggedright\arraybackslash}W{.23}>{\raggedright\arraybackslash}W{.17}@{}}
\rowcolor{brand}\ptHead{\VAR{L.id}} & \ptHead{\VAR{L.title}} & \ptHead{\VAR{L.severity}} & \ptHead{\VAR{L.retest_result}} & \ptHead{\VAR{L.retest_date}} \\
\endhead
\BLOCK{for f in findings}
\hyperref[f:\VAR{f.id}]{\small\texttt{\VAR{f.display_id}}} & \raggedright\small \VAR{f.title} & \ptSeverity{\VAR{f.severity}}{\VAR{sev_labels[f.severity]}} & \ptRetest{\VAR{(f.retest_result or 'none')|replace('_', '')}}{\VAR{retest_labels[f.retest_result]}} & \small\texttt{\VAR{f.retest_date or ''}} \\
\BLOCK{endfor}
\end{longtable}
\rowcolors{1}{}{}
\BLOCK{else}
\VAR{L.no_findings}
\BLOCK{endif}
\BLOCK{endmacro}

\BLOCK{macro charts()}
\BLOCK{set total = findings|length}
%# --- findings by severity: donut + legend
\begin{tcolorbox}[enhanced, colback=white, colframe=line, boxrule=0.5pt, arc=3pt, title={\VAR{L.chart_severity}},
  coltitle=ink, colbacktitle=white, fonttitle=\small\bfseries, titlerule=0.5pt, titlerule style={line}]
\centering
\begin{tikzpicture}[x=1mm, y=1mm]
\BLOCK{if total}
\BLOCK{set ns = namespace(angle=90.0)}
\BLOCK{for s in severities}\BLOCK{if counts[s]}
\BLOCK{set sweep = 360.0 * counts[s] / total}
  \fill[sev\VAR{s}] (0,0) -- (\VAR{ns.angle|round(3)}:22) arc[start angle=\VAR{ns.angle|round(3)}, delta angle=\VAR{(-sweep)|round(3)}, radius=22] -- cycle;
\BLOCK{set ns.angle = ns.angle - sweep}
\BLOCK{endif}\BLOCK{endfor}
\BLOCK{else}
  \fill[soft] (0,0) circle (22);
\BLOCK{endif}
  \fill[white] (0,0) circle (12);
  \node[font=\Large\bfseries] at (0, 1.5) {\VAR{total}};
  \node[font=\scriptsize, text=muted] at (0, -3.5) {\VAR{L.number_of_findings}};
\BLOCK{for s in severities}
  \fill[sev\VAR{s}, rounded corners=0.6pt] (40, \VAR{(15 - 7 * loop.index0 - 1.8)|round(2)}) rectangle (43.6, \VAR{(15 - 7 * loop.index0 + 1.8)|round(2)});
  \node[anchor=west, font=\small] at (46, \VAR{(15 - 7 * loop.index0)|round(2)}) {\VAR{sev_labels[s]}};
  \node[anchor=east, font=\small\bfseries] at (92, \VAR{(15 - 7 * loop.index0)|round(2)}) {\VAR{counts[s]}};
\BLOCK{endfor}
\end{tikzpicture}
\end{tcolorbox}

\BLOCK{if per_target}
%# --- findings per target: stacked columns
\BLOCK{set tmax = ([1] + per_target|map(attribute='total')|list)|max}
\BLOCK{set colw = [10, (120 / per_target|length) * 0.55]|min}
\BLOCK{set step = 120 / per_target|length}
\begin{tcolorbox}[enhanced, colback=white, colframe=line, boxrule=0.5pt, arc=3pt, title={\VAR{L.chart_targets}},
  coltitle=ink, colbacktitle=white, fonttitle=\small\bfseries, titlerule=0.5pt, titlerule style={line}]
\centering
\begin{tikzpicture}[x=1mm, y=1mm]
\BLOCK{set kstep = 1 if tmax <= 10 else (tmax + 9) // 10}
\BLOCK{for k in range(0, tmax + 1, kstep)}
  \draw[line] (0, \VAR{(40 * k / tmax)|round(2)}) -- (122, \VAR{(40 * k / tmax)|round(2)});
  \node[anchor=east, font=\scriptsize, text=muted] at (-1, \VAR{(40 * k / tmax)|round(2)}) {\VAR{k}};
\BLOCK{endfor}
\BLOCK{for t in per_target}
\BLOCK{set x0 = step * loop.index0 + (step - colw) / 2 + 1}
\BLOCK{set ns = namespace(y=0.0)}
\BLOCK{for s in severities}\BLOCK{if t.counts[s]}
\BLOCK{set h = 40 * t.counts[s] / tmax}
  \fill[sev\VAR{s}] (\VAR{x0|round(2)}, \VAR{ns.y|round(2)}) rectangle (\VAR{(x0 + colw)|round(2)}, \VAR{(ns.y + h)|round(2)});
\BLOCK{set ns.y = ns.y + h}
\BLOCK{endif}\BLOCK{endfor}
  \node[anchor=north, font=\scriptsize, text width=\VAR{(step - 2)|round(2)}mm, align=center] at (\VAR{(x0 + colw / 2)|round(2)}, -1.5) {\VAR{t.name}};
\BLOCK{endfor}
\end{tikzpicture}

\smallskip
{\footnotesize\BLOCK{for s in severities}\textcolor{sev\VAR{s}}{\rule{2.4mm}{2.4mm}}\,\VAR{sev_labels[s]}\quad\BLOCK{endfor}}
\end{tcolorbox}
\BLOCK{endif}

%# --- CVSS rating versus severity: bubble matrix
\begin{tcolorbox}[enhanced, colback=white, colframe=line, boxrule=0.5pt, arc=3pt, title={\VAR{L.chart_matrix}},
  coltitle=ink, colbacktitle=white, fonttitle=\small\bfseries, titlerule=0.5pt, titlerule style={line}]
\centering
\begin{tikzpicture}[x=1mm, y=1mm]
\BLOCK{set cols = ['none', 'low', 'medium', 'high', 'critical', '']}
\BLOCK{for s in severities}
  \node[anchor=east, font=\small] at (24, \VAR{(-9 * loop.index0)|round(2)}) {\VAR{sev_labels[s]}};
\BLOCK{set row = loop.index0}
\BLOCK{for r in cols}
\BLOCK{set n = matrix[s][r]}
\BLOCK{set cx = 33 + 19 * loop.index0}
  \draw[line, fill=soft] (\VAR{(cx - 9)|round(2)}, \VAR{(-9 * row - 4.2)|round(2)}) rectangle (\VAR{(cx + 9)|round(2)}, \VAR{(-9 * row + 4.2)|round(2)});
\BLOCK{if n}
  \fill[sev\VAR{s}] (\VAR{cx}, \VAR{(-9 * row)|round(2)}) circle (\VAR{(1.6 + 2.2 * (n / [total, 1]|max) ** 0.5)|round(2)});
  \node[font=\scriptsize\bfseries, text=white] at (\VAR{cx}, \VAR{(-9 * row)|round(2)}) {\VAR{n}};
\BLOCK{endif}
\BLOCK{endfor}
\BLOCK{endfor}
\BLOCK{for r in cols}
  \node[font=\scriptsize, text=muted, align=center, text width=17mm] at (\VAR{33 + 19 * loop.index0}, -42) {\BLOCK{if r}\VAR{r|rating_label(lang)}\BLOCK{else}\VAR{L.no_cvss}\BLOCK{endif}};
\BLOCK{endfor}
  \node[font=\scriptsize\bfseries, text=muted] at (80.5, -47.5) {\VAR{L.cvss_score}};
  \node[font=\scriptsize\bfseries, text=muted, rotate=90] at (-4, -18) {\VAR{L.severity}};
\end{tikzpicture}
\end{tcolorbox}
\BLOCK{endmacro}

\BLOCK{macro finding(f, grouped=False, first=False)}
\BLOCK{if not first}\clearpage\BLOCK{endif}
\phantomsection\label{f:\VAR{f.id}}
\BLOCK{if grouped}
\subsubsection{\texorpdfstring{\texttt{\VAR{f.display_id}} \VAR{f.title}\BLOCK{if f.status != 'open'}\enspace{\normalfont\small\color{muted}\MakeUppercase{\VAR{status_labels[f.status]}}}\BLOCK{endif}}{\VAR{f.display_id} \VAR{f.title}}}
\BLOCK{else}
\addcontentsline{toc}{subsection}{\VAR{f.display_id}\quad \VAR{f.title}}
{\small\color{muted}\texttt{\bfseries{}\VAR{f.display_id}}}\par\vspace{2pt}
{\Large\bfseries \VAR{f.title}\par}\vspace{8pt}
\BLOCK{endif}

%# Risk band: the severity in large type on its colour, the scores beside it.
\begin{tcolorbox}[enhanced, colback=white, colframe=line, boxrule=0.5pt, arc=3pt, left=8pt, right=8pt, top=3pt, bottom=3pt,
  title={\ptRiskBand{\VAR{sev_labels[f.severity]}}{\BLOCK{if f.cvss_score is not none}CVSS \VAR{f.cvss_version} \textbf{\VAR{f.cvss_score|score}}\BLOCK{endif}\BLOCK{if f.cwe}\quad \VAR{f.cwe}\BLOCK{endif}\quad \VAR{status_labels[f.status]}}},
  colbacktitle=sev\VAR{f.severity}, coltitle=white, toptitle=6pt, bottomtitle=6pt, lefttitle=10pt, righttitle=10pt, titlerule=0pt]
\begin{tabularx}{\linewidth}{@{}>{\leavevmode\color{muted}\small\bfseries}p{34mm}L@{}}
\BLOCK{if f.targets}
\VAR{L.affected} & \BLOCK{for t in f.targets}\VAR{t.name}\BLOCK{if t.address} {\small\color{muted}\texttt{\VAR{t.address}}}\BLOCK{endif}\BLOCK{if not loop.last}\newline \BLOCK{endif}\BLOCK{endfor} \\
\BLOCK{endif}
\BLOCK{if f.cvss_vector}
CVSS \VAR{f.cvss_version} & \ptCvss{\VAR{f|cvss_rating_key}}{\VAR{f.cvss_score|score}}{\VAR{f|cvss_rating_label(lang)}}\quad{\footnotesize\texttt{\VAR{f.cvss_vector}}}\newline
\BLOCK{for g in f.cvss_vector|vector_strip}\BLOCK{for c in g.cells}\ptCell{\VAR{c.tone}}{\VAR{c.key}}{\VAR{c.value}}\BLOCK{endfor}\hspace{4pt}\BLOCK{endfor} \\
\BLOCK{endif}
\BLOCK{if f.cwe and not grouped}
\VAR{L.cwe} & \VAR{f.cwe|url_cwe}\BLOCK{if f.cwe|cwe_name}: \VAR{f.cwe|cwe_name}\BLOCK{endif} \\
\BLOCK{endif}
\VAR{L.status} & \VAR{status_labels[f.status]}\BLOCK{if f.confidence}\qquad{\color{muted}\small\bfseries \VAR{L.confidence}:} \VAR{conf_labels[f.confidence]}\BLOCK{endif} \\
\BLOCK{if f.retest_result}
\VAR{L.retest} & \ptRetest{\VAR{f.retest_result|replace('_', '')}}{\VAR{retest_labels[f.retest_result]}}\BLOCK{if f.retest_date}\quad{\small\texttt{\VAR{f.retest_date}}}\BLOCK{endif} \\
\BLOCK{endif}
\end{tabularx}
\end{tcolorbox}

\BLOCK{if f.retest_notes}
\ptBlockTitle{\VAR{L.retest_notes}}
\VAR{f.retest_notes|richtext}
\BLOCK{endif}

\BLOCK{for key in ['description', 'impact', 'technical_details', 'poc']}
\BLOCK{if f[key]}
\ptBlockTitle{\VAR{L[key]}}
\VAR{f[key]|richtext}
\BLOCK{endif}
\BLOCK{endfor}

\BLOCK{for a in f.attachments}
\begin{center}
\VAR{a|image}
\captionof{figure}{\VAR{a.caption}}
\end{center}
\BLOCK{endfor}

\BLOCK{if f.evidence}
\ptBlockTitle{\VAR{L.evidence}}
\BLOCK{for e in f.evidence}
\needspace{6\baselineskip}
\BLOCK{if e.title}{\small\bfseries \VAR{e.title}}\par\vspace{3pt}\BLOCK{endif}
\ptHttp{\VAR{L.request}}{\VAR{e.request|verbatim_file('request')}}
\BLOCK{if e.response}
\ptHttp{\VAR{L.response}}{\VAR{e.response|verbatim_file('response')}}
\BLOCK{endif}
\BLOCK{endfor}
\BLOCK{endif}

%# Evidence files: name, size and SHA-256 (in two halves, so it always fits its column) for integrity checks.
%# One fixed-width tabular per row, not a longtable: a page can break between any two rows. (A longtable after a
%# \ptBlockTitle on the same page can make TeX break at that title's \needspace and repeat the table header there.)
%# \leavevmode: a \color at the very start of a p-cell would otherwise make the header row taller.
\BLOCK{if f.files}
\ptBlockTitle{\VAR{L.attached_files}}
\begingroup\centering\setlength{\parskip}{0pt}
\begin{tabular}{@{}W{.44}>{\raggedleft\arraybackslash}W{.12}W{.44}@{}}
\rowcolor{brand}\leavevmode\ptHead{\VAR{L.file_name}} & \leavevmode\ptHead{\VAR{L.file_size}} & \leavevmode\ptHead{\VAR{L.sha256}} \\
\end{tabular}\par\nopagebreak
\BLOCK{for x in f.files}
\nointerlineskip\begin{tabular}{@{}W{.44}>{\raggedleft\arraybackslash}W{.12}W{.44}@{}}
\BLOCK{if loop.index is even}\rowcolor{soft}\BLOCK{endif}\raggedright\small \VAR{x.filename|breakable}\BLOCK{if x.description}\newline{\footnotesize\color{muted}\VAR{x.description}}\BLOCK{endif} & \small \VAR{x.size|filesize} & \footnotesize\texttt{\VAR{x.sha256[:32]}}\newline\texttt{\VAR{x.sha256[32:]}} \\
\end{tabular}\par
\BLOCK{endfor}
\par\endgroup
\BLOCK{endif}

\BLOCK{if f.remediation}
\ptBlockTitle{\VAR{L.remediation}}
\VAR{f.remediation|richtext}
\BLOCK{endif}

\BLOCK{if f.refs}
\ptBlockTitle{\VAR{L.references}}
\VAR{f.refs|reference_list}
\BLOCK{endif}

\BLOCK{if show_finding_history and f.history}
\ptBlockTitle{\VAR{L.finding_history}}
\rowcolors{1}{soft}{white}
\begin{tabularx}{\linewidth}{@{}p{24mm}\BLOCK{if f.history_has_author}p{38mm}\BLOCK{endif}L@{}}
\BLOCK{for h in f.history}
\texttt{\small\VAR{h.date or ''}} & \BLOCK{if f.history_has_author}\small\VAR{h.author} & \BLOCK{endif}\small\VAR{h.text} \\
\BLOCK{endfor}
\end{tabularx}
\rowcolors{1}{}{}
\BLOCK{endif}
\BLOCK{endmacro}

\BLOCK{macro findings_body()}
\BLOCK{if not findings}
\VAR{L.no_findings}
\BLOCK{elif groups}
\BLOCK{for g in groups}
\BLOCK{if not loop.first}\clearpage\BLOCK{endif}
\subsection{\VAR{g.title}}
\VAR{g.description|richtext}
\BLOCK{for f in g.findings}
\VAR{finding(f, grouped=True, first=loop.first)}
\BLOCK{endfor}
\BLOCK{if g.mitigation}
\subsubsection{\VAR{L.mitigation}}
\VAR{g.mitigation|richtext}
\BLOCK{endif}
\BLOCK{endfor}
\BLOCK{else}
\BLOCK{for f in findings}
\VAR{finding(f, first=loop.first)}
\BLOCK{endfor}
\BLOCK{endif}
\BLOCK{endmacro}

%# ========================================================================================== body ==
%# The order and titles come from the report structure (Report -> Structure in the app).
\BLOCK{for sec in sections}
\VAR{heading(sec)}
\BLOCK{if sec.kind == 'executive_summary'}
\VAR{report.executive_summary|richtext(sec.numbered)}
\BLOCK{elif sec.kind == 'scope'}
\VAR{report.scope|richtext(sec.numbered)}
\BLOCK{if targets}
\subsection*{\VAR{L.targets}}
\VAR{targets_table()}
\BLOCK{endif}
\BLOCK{elif sec.kind == 'methodology'}
\VAR{report.methodology|richtext(sec.numbered)}
\BLOCK{elif sec.kind == 'summary'}
\VAR{severity_bars()}
\vspace{0.6em}
\VAR{findings_table()}
\BLOCK{elif sec.kind == 'findings_table'}
\VAR{findings_table()}
\BLOCK{elif sec.kind == 'charts'}
\VAR{charts()}
\BLOCK{elif sec.kind == 'retest_summary'}
\VAR{retest_summary()}
\BLOCK{elif sec.kind == 'findings'}
\VAR{findings_body()}
\BLOCK{endif}
\VAR{sec.body|richtext(sec.numbered)}
\BLOCK{endfor}

\end{document}
